\makeabstract
{
Galaxy Clusters as a Bridge between Theory and Observation
}
{
Ryogo Katahira (1), Tomomi Sunayama (2)
}
{
\textsuperscript{1}Amherst College\\
\textsuperscript{2}Academia Sinica Institute of Astronomy and Astrophysics
}
{
In current cosmology, ΛCDM is considered as a standard model of cosmology that best explains the evolution of the universe with only six parameters. Dark energy (Λ) and cold dark matter (CDM) lie at the core of ΛCDM, yet their identities still remain as a question. Constraining the cosmological parameters is crucial to understand the nature of dark energy, and galaxy clusters, which are formed from initial density fluctuations in the early universe, are considered as  one of the most powerful cosmological probes. However, using clusters as a cosmological probe requires mass-richness relation (P(λ|M)), since the cluster mass is not a direct observable. In this research, we build a pipeline that constrain parameters of P(λ|M) and seek ways to improve in a framework called full forward modeling.
}
{
In order to build a pipeline, we used a mock catalog of clusters from Sunayama et al. 2020. We used three observables to constrain the parameters of P(λ|M): abundance , stacked lensing signal, and clustering. On the other hand, we computed the observables from cosmological parameters and P(λ|M) using dark\_emulator, a cosmological emulator developed in Nishimachi et al. 2019. To evaluate the parameters, we apply Bayesian analysis through the Markov chain Monte Carlo (MCMC), which collects “likely” parameter values by random-walking and returns a probability distribution of parameter values called posterior.
}
{
From our best run of MCMC, most parameters are normally distributed and constrained well with small 1σ, yet one parameter shows two peaks in its probability distribution.
When comparing P(λ|M), the predicted P(λ|M) is narrower in every λ bin. Potential causes include incorrectly assuming the same relationship across all λ and mass bias from weak lensing.
As future steps, we will seek ways to eliminate the slight disagreements with the sample and would like to add cosmological parameters as free parameters. We will then run the pipeline with the CAMIRA cluster catalog, possibly with multiple redshift values.
}
{
In this research, we build a pipeline that constrain parameters of P(λ|M) for future analysis with real cluster catalog and cosmological parameters. The results of MCMC mostly match the sample with a slight disagreement in width of posterior distribution. The parameter values also show slight disagreement from previous studies, yet the uncertainties became smaller with the presence of clustering in addition to abundance and lensing signal. In the future, we will run the pipeline with real data and incorporate cosmological parameters as free parameters.
}

\newpage
